\documentclass[12pt,a4paper]{article}
\setlength{\parindent}{0pt}

\usepackage{graphicx}
\usepackage{amsmath}
\usepackage{amssymb}
\usepackage{geometry}
\usepackage{hyperref}
\usepackage{fancyhdr}
\usepackage{titlesec}
\usepackage{booktabs}
\usepackage{array}
\usepackage{bm}
\usepackage{float}
\usepackage{graphicx}


\geometry{left=3cm,right=3cm,top=3cm,bottom=3cm}


\pagestyle{fancy}
\fancyhf{}
\fancyhead[L]{copf2}
\fancyhead[R]{Sheng'en Li}
\fancyfoot[C]{\thepage}

\titleformat{\section}{\large\bfseries}{\thesection}{1em}{}
\titleformat{\subsection}{\normalsize\bfseries}{\thesubsection}{1em}{}

\begin{document}
% =========================
%  Skeleton: COPF core + heavy-tail extensions (SIAM-theory oriented)
% =========================

% =========================
% Theorem C blueprint (SIAM-theory oriented)
% Random graph process -> (kappa, tau_mix, heavy-tail moments)
% =========================

\section{Random-Graph Foundations for Bounded Dependence}
\label{sec:random_graph_foundations}

\subsection{COPF dependence objects (recap)}
Following COPF, we index an audited instance by
$u=(i,j,t)$ (a candidate dyad at time $t$), and assume
an $(r_{\mathrm{int}},L_{\mathrm{int}})$-local interference neighborhood $I(u)$
that is \emph{computable} from $G_{\le t}$ and the candidate rule.
A (random) dependency graph is defined on instances within a window:
connect $u,u'$ if $u'\in I(u)$ or $u\in I(u')$.
Let $\kappa$ denote its maximum degree.

We also use a temporal mixing notion (e.g., $\beta$-mixing, or function-class mixing on a
finite test family) with mixing time $\tau_{\mathrm{mix}}$.

\subsection{A dynamic random graph generative model}
\label{subsec:dyn_random_graph}

We consider an evolving graph process $\{G_t\}_{t\ge 0}$ on vertex set $V_t$
(possibly growing with $t$). The model is left modular:

\begin{assumption}[Dynamic random graph model $\mathsf{GRM}(\Theta)$]
\label{ass:grm}
The evolution $\{G_t\}$ is generated by a parameterized random process $\mathsf{GRM}(\Theta)$
(e.g., preferential attachment with choice, spatial random graphs, or a sparse inhomogeneous random graph).
The model admits:
\begin{enumerate}
\item \textbf{Local volume growth control.} For each fixed radius $r$,
there exists a deterministic function $b_r(n)$ such that with probability at least $1-\delta$,
\[
\max_{t\le T}\ \max_{v\in V_t} \ |B_r(v;G_t)| \ \le\ b_r(n_T),
\]
where $n_T:=|V_T|$ and $B_r(v;G_t)$ is the radius-$r$ ball in $G_t$.
\item \textbf{(Optional) degree tail exponent.} The marginal degree distribution obeys a power-law tail
with exponent $\gamma>1$ in the sense that, for large $x$,
$\Pr(\deg(v)\ge x)\approx x^{-(\gamma-1)}$ (or an appropriate high-probability surrogate).
\item \textbf{(Optional) fast mixing of the driving chain.} The event process is generated by
an underlying Markov chain with spectral gap $\lambda>0$ (defined precisely below), yielding
temporal dependence decay.
\end{enumerate}
\end{assumption}

% FanWei-core: pick a concrete GRM (PA-with-choice / spatial) and prove b_r(n) polylog,
% plus (optional) tail exponent gamma and a usable mixing statement.


\subsection{Candidate rule + Top-K policy (policy layer)}
We allow candidate sets that depend on the evolving graph.
At time $t$, given a source/anchor $u_t$, define the candidate set $C_t=C_t(u_t,G_{\le t})$.
Two canonical special cases:
(i) bounded-budget negatives (COPF-style), $|C_t|\le N+1$,
(ii) neighborhood candidates, $C_t \subseteq B_{r_{\mathrm{cand}}}(u_t;G_{\le t})$.

We consider a stochastic Top-$K$ exposure policy:
with probability $\epsilon$ we sample a size-$K$ slate uniformly without replacement from $C_t$;
otherwise we sample a Top-$K$ slate from a score-based distribution (e.g., Plackett--Luce).
We log marginal entry propensities $e_t(v)=\Pr(v\in D_t\mid \mathcal{H}_t)$ for all $v\in C_t$.

% NOTE: COPF clips propensities for numerical stability; here we analyze the unclipped regime.

\subsection{A reduction lemma: graph volume growth $\Rightarrow$ bounded dependency degree}
\label{subsec:kappa_reduction}

\begin{lemma}[Design-to-structure reduction for $\kappa$]
\label{lem:kappa_reduction}
Assume $(r_{\mathrm{int}},L_{\mathrm{int}})$-local interference neighborhoods $I(u)$
only include instances whose endpoints lie within $r_{\mathrm{int}}$ hops of $(i,j)$
and within the last $L_{\mathrm{int}}$ rounds (as in COPF).
Let $M_t:=|C_t|$ and $M:=\max_{t\le T} M_t$.
Then, for any window of length $W$,
the dependency-graph maximum degree satisfies
\[
\kappa \ \le\ c \cdot (2L_{\mathrm{int}}+1)\cdot M \cdot
\max_{t\le T}\max_{v\in V_t}\ |B_{2r_{\mathrm{int}}}(v;G_t)|,
\]
for a universal constant $c$ depending only on the precise instance indexing
(e.g., whether both endpoints contribute symmetrically).
\end{lemma}

% FanWei-core: sharpen constants / tighten the factor 2r_int,
% and upgrade the bound to chromatic number / degeneracy if that gives better concentration.

\subsection{Theorem C.1: sparse interference (high probability)}
\label{subsec:C1}

\begin{theorem}[C.1: Sparse interference / bounded $\kappa$ under $\mathsf{GRM}(\Theta)$]
\label{thm:C1_sparse_interference}
Under Assumption~\ref{ass:grm} (local volume growth control) and Lemma~\ref{lem:kappa_reduction},
for any $\delta\in(0,1)$, with probability at least $1-\delta$,
\[
\kappa \ \le\ c\cdot(2L_{\mathrm{int}}+1)\cdot M \cdot b_{2r_{\mathrm{int}}}(n_T).
\]
In particular, if $\mathsf{GRM}(\Theta)$ yields $b_{2r_{\mathrm{int}}}(n)=\polylog(n)$,
then $\kappa=\polylog(n_T)$ with high probability even when $G_t$ has heavy-tailed degrees.
\end{theorem}

% FanWei-core: instantiate GRM so that b_r(n) = O((\log n)^c) w.h.p.
% e.g. PA-with-choice / spatial graphs often have branching-process approximations for fixed-radius balls.


\subsection{Theorem C.2: mixing time via spectral gap}
\label{subsec:C2}

We model temporal dependence through a Markov driving chain on a state space $\mathcal{S}$
that generates the logged tuples $Z_t$ (contexts, actions, outcomes, etc.).
Let $P$ be its transition kernel and $\lambda$ its spectral gap.

\begin{theorem}[C.2: Temporal mixing from spectral gap]
\label{thm:C2_mixing_spectral}
Assume the driving chain is reversible with spectral gap $\lambda>0$ and stationary distribution $\pi$.
Then its $\beta$-mixing coefficients decay geometrically:
\[
\beta(\ell)\ \le\ C \cdot (1-\lambda)^\ell,
\]
for a universal constant $C$, and thus the mixing time satisfies
\[
\tau_{\mathrm{mix}}(\varepsilon)\ \le\ \frac{1}{\lambda}\log\!\Bigl(\frac{C}{\varepsilon}\Bigr).
\]
If additionally the induced interference/dependency structure is sparse as in Theorem~\ref{thm:C1_sparse_interference},
one can combine the dependence-through-time (via $\tau_{\mathrm{mix}}$) and dependence-within-window
(via $\kappa$) to obtain explicit end-to-end concentration radii.
\end{theorem}

% FanWei-core: for a concrete GRM + event dynamics, prove a nontrivial lower bound on lambda,
% ideally uniform in n (or at least polylog rates). This is where spectral graph theory shines.


\subsection{Theorem C.3: propensity moments and heavy-tail index}
\label{subsec:C3}

This result links graph heavy tails to inverse-propensity heavy tails.

\begin{theorem}[C.3: Inverse propensity moments under neighborhood candidates]
\label{thm:C3_propensity_moments}
Assume neighborhood candidate sets: $C_t \subseteq B_{r_{\mathrm{cand}}}(u_t;G_{\le t})$
and the exposure policy has an $\epsilon$-uniform component that samples a size-$K$ slate
uniformly without replacement from $C_t$.
Then for any $v\in C_t$,
\[
e_t(v)\ \ge\ \epsilon\cdot \frac{K}{|C_t|}.
\]
Consequently,
\[
\Bigl(\frac{1}{e_t(v)}\Bigr)^p \ \le\ \Bigl(\frac{|C_t|}{\epsilon K}\Bigr)^p.
\]
If $|C_t|$ is controlled by local graph volume, e.g. $|C_t|\le |B_{r_{\mathrm{cand}}}(u_t;G_{\le t})|$,
then
\[
\E\Bigl[\Bigl(\frac{1}{e_t(v)}\Bigr)^p\Bigr]\ <\ \infty
\quad \text{whenever}\quad
\E\bigl[|B_{r_{\mathrm{cand}}}(u_t;G_{\le t})|^p\bigr]\ <\ \infty.
\]
In particular, under a power-law degree tail with exponent $\gamma$,
a sufficient condition is $p<\gamma-1$ (yielding a phase transition in the highest finite moment).
\end{theorem}

\begin{remark}[Phase transition viewpoint]
Theorem~\ref{thm:C3_propensity_moments} implies a critical moment threshold:
as the graph becomes more heavy-tailed (smaller $\gamma$), fewer moments of $1/e_t$ remain finite.
This directly determines the admissible heavy-tail index $p\in(1,2]$ for robust estimation/optimization.
\end{remark}

% FanWei-core: make the p<gamma-1 condition sharp for the chosen GRM,
% and refine the bound beyond the uniform component (handle score-based PL Top-K inclusion).

% TODO (FanWei, central): pick a clean structural family where both kappa and mixing are provably bounded,
% e.g. bounded degeneracy + limited-radius neighborhoods; or a random graph model where dependence decays.

% =========================
% End skeleton
% =========================

\end{document}
